\section{The Formal Grammar: Bounded Executive Action with Typed Rollback Witnesses}
\label{sec:grammar}

\subsection{The construction}

The grammar this Article proposes formalizes a four-component requirement for any delegated emergency authority.
At the moment an authority is exercised, the exercising actor must specify four things:
the substantive action to be performed,
a duration after which the authority lapses,
a witness that the prior state remains constructively recoverable,
and a predicate naming the actors whose concurrence is required to admit the action into the polity's history.
The four components correspond to the doctrinal categories of substantive action, sunset, reversibility,
and authorization procedure as those categories appear in existing constitutional and administrative law.
The grammar's contribution is not the invention of any of the four categories,
each of which has a long doctrinal pedigree.
The contribution is the requirement that all four be present, simultaneously and at the moment of admission,
as conditions of an exercise of authority that the legal order will recognize.

The framing is structural rather than substantive.
The grammar does not specify what content any of the four components should have in any particular case.
The substantive action, the duration, the form of the rollback witness, and the membership of the quorum
are all matters of democratic political choice that the grammar takes as exogenous.
What the grammar adds is the requirement that, whatever content the political process assigns to the four
components, the assignment must occur at the moment of admission and must be type-checkable at that moment
rather than deferred to a later audit, review, or lapse step.
The first three components correspond to features for which the existing statutory record provides
systematic instances; the fourth, the quorum predicate, is the component that captures the
Ackerman-Gersen structural-correction proposals on supermajority defaults and duration clauses for emergency-power continuation.

With that legal framing in place, the formal expression of the grammar can be stated compactly.
A bounded executive action is a four-tuple
\[
a = (\text{act}, \TTL, w, q)
\]
where $\text{act}$ is the action to be performed, $\TTL$ is the time-to-live bounding the action's
period of effect, $w$ is a typed rollback witness, and $q$ is the quorum predicate specifying which actors
must concur in the action's admission.
The action $a$ is admissible into the polity's history only if the four-tuple is well-formed in the sense
defined below.

The typed rollback witness $w$ is the load-bearing component of the construction.
A typed rollback witness for an action $\text{act}$ is a proof, constructible at the moment of admission of $a$,
that the polity's state after the execution of $\text{act}$ admits a transformation back to its state before
the execution of $\text{act}$,
and that the inverse transformation is itself an action the polity is authorized to perform.
The witness is not a promise to roll back, nor a procedure by which a later actor may request rollback.
The witness is a proof that the rollback is constructible.

The distinction between a promise and a proof matters.
A promise is a statement about future action that is binding only insofar as the political process
that enforces it operates as designed.
A proof is a constructive object that, by its existence, demonstrates that the future action is available.
A promise can be broken; a proof either exists or does not.
An action admitted without a proof of its rollback's constructibility may turn out, at the moment rollback is required,
to be irreversible.
An action admitted with a proof of its rollback's constructibility is, by the construction of the proof,
reversible by available means.

The grammar makes the unconstructive case unconstructible.
An action $a$ for which no rollback witness $w$ exists cannot be constructed as a bounded executive action.
The exercise of an authority granting such an action fails to type-check.
The failure is not at the audit step, after the action has been taken,
but at the admission step, before the action has been taken.
The typed-rollback discipline applies at construction time, not at execution time.

\subsection{The discipline applied to the case studies}

Each of the five case studies in Part~\ref{sec:cases} is, on the grammar's framing,
a case in which one or more components of the four-tuple is absent from the statutory record.

Weimar Article~48 supplied the substantive component ($\text{act}$), in the form of the authority to take
``necessary measures.'' The article did not supply a typed rollback witness ($w$),
and the time-to-live ($\TTL$) was implicit and politically determined rather than typed.
The quorum predicate ($q$) was the Reichstag's authority to demand suspension,
which was a post-hoc check rather than a precondition of admission.
The grammar's reading of the article is that the action component was present but the other three components
were either absent or unenforced at the type level.

The DMCA Section~512 notice-and-takedown regime, operating alongside Section~230's liability shield
for platform self-regulation, supplies $\text{act}$ (the removal) and an implicit $\TTL$
(the action takes effect immediately and remains in effect until or unless restored).
The witness $w$ is not required at construction time.
The Section~512(g) counter-notice procedure is a post-hoc check, not a precondition of admission.
The quorum $q$ is a single private actor, the platform's content-moderation function.

GDPR Article~17 supplies $\text{act}$ and a procedural $q$ (the administrative process by which an erasure order
is issued).
The witness $w$ is, in practice, not constructible at the moment of admission because the erasure is, in the
present technical and operational regime, irreversible by available means.
The grammar's reading is that the construction of an Article~17 erasure order, on the present technical regime,
fails the typed-rollback test by default.

FISA Section~702 supplies $\text{act}$ (the surveillance authorization) and an implicit $q$ (Attorney General
plus Director of National Intelligence).
The witness $w$ is absent because the surveillance, once executed, produces a state of affairs --
information in the holdings of the collecting agency -- for which no constructive rollback is statutorily
required and, on the present technical regime, no constructive rollback is operationally available.

The 2001 AUMF supplies $\text{act}$ (use of military force) and a $q$ (presidential determination of the
target's covered status).
The witness $w$ is absent.
The $\TTL$ is absent in the statutory text.
The grammar's reading of the AUMF is that three of the four components are either absent or absent at the
type level, with only the action component present in statutory form.

The grammar's prescription, in each case, is structural rather than substantive.
The grammar does not specify what the substantive action should be, what the time-to-live should be,
what the rollback witness should look like, or what the quorum should consist of.
It specifies only that all four components must be present, constructible, and type-checked at the moment
of admission.
The substantive content of each component is a matter of democratic political determination,
which the grammar deliberately leaves to the legislative process that enacts the underlying authority.

\subsection{What the grammar does and does not claim}

The grammar makes a structural claim.
It does not make the substantive claim that any particular delegated emergency authority is desirable or undesirable.
It does not specify which actions should be permissible under emergency authority,
which time-to-live values are appropriate, what content a rollback witness should have, or which actors should
be authorized to concur.
These are questions of democratic political choice that the grammar takes as exogenous.

The grammar also does not claim that the construction is operationally costless.
Requiring the construction of a typed rollback witness at the moment of admission imposes a structural cost
on the actor seeking admission.
In some cases the cost will be small -- a platform that maintains versioned content storage can construct
a rollback witness for a removal as a routine matter.
In other cases the cost will be large -- a surveillance authority whose technical implementation does not
permit constructive reversal of collection may not be able to exhibit a witness without substantial re-engineering
of the underlying collection systems.
The cost is a real consideration that any application of the grammar must address.

What the grammar claims is the narrower observation that the cost is currently being shifted onto the polity
rather than borne by the actor seeking admission.
When an emergency authority is exercised without a typed rollback witness,
the cost of the resulting irreversibility falls on the polity, not on the actor.
The polity bears the cost in the form of permanent alterations to the state of affairs that the authority was
intended to alter only temporarily.
The grammar shifts the cost from the polity to the actor.
The actor must, at admission, either construct the witness or refrain from seeking admission.
The shift is, on the Ackerman-Gersen structural-correction framing, the appropriate distribution of the cost,
because the actor seeking admission is the actor in the best position to bear it.

The grammar's relationship to the Schmittian objection -- that the content of the exception is not formalizable --
is that the grammar formalizes only the wrapper.
The substantive content of the action, the substantive determination of whether an exception obtains,
the substantive question of what state of affairs the polity should return to: these are political determinations
that the grammar takes as inputs.
What the grammar provides is a structural skeleton within which those determinations are made.
The skeleton is constructible even if the content is not.
A Schmittian critic of the grammar would observe, correctly, that the most consequential element of an emergency
authority is the determination that the exception obtains.
The Article's response is that the determination is preserved as a political determination in the grammar;
the grammar adds the requirement that the determination be accompanied by a typed witness for the path back.
The witness does not prejudge the determination.
The witness conditions only the form of the action that follows from the determination.
